% theme: conversions_among_moles_mass_gas_volume_and_particles
\MajorQuestion{物質量・質量・気体の体積・粒子数の換算}
\SetRowSpace{3mm}

\TextTall{原子量は ${\rm H}=1.0$、${\rm C}=12$、${\rm N}=14$、${\rm O}=16$、${\rm Na}=23$、${\rm Cl}=35.5$ とする。また、アボガドロ定数を $6.0\times10^{23}\,{\rm mol}^{-1}$ とし、標準状態における気体 $1\,{\rm mol}$ の体積を $22.4\,{\rm L}$ とする。}

\TypeQuestion{次の物質量を質量に換算しなさい。}{mol_to_mass}
\QQRow{酸素 $0.50\,{\rm mol}$}{$16\,{\rm g}$}{二酸化炭素 $1.5\,{\rm mol}$}{$66\,{\rm g}$}
\QQRow{塩化ナトリウム $0.20\,{\rm mol}$}{$11.7\,{\rm g}$}{水 $2.0\,{\rm mol}$}{$36\,{\rm g}$}
\QQRow{メタン $0.75\,{\rm mol}$}{$12\,{\rm g}$}{アンモニア $0.40\,{\rm mol}$}{$6.8\,{\rm g}$}

\TypeQuestion{次の質量を物質量に換算しなさい。}{mass_to_mol}
\QQRow{水 $9.0\,{\rm g}$}{$0.50\,{\rm mol}$}{アンモニア $34\,{\rm g}$}{$2.0\,{\rm mol}$}
\QQRow{窒素 $7.0\,{\rm g}$}{$0.25\,{\rm mol}$}{塩化ナトリウム $23.4\,{\rm g}$}{$0.40\,{\rm mol}$}
\QQRow{二酸化炭素 $13.2\,{\rm g}$}{$0.30\,{\rm mol}$}{酸素 $48\,{\rm g}$}{$1.5\,{\rm mol}$}

\TypeQuestion{次の物質量の気体が標準状態で占める体積を求めなさい。}{mol_to_gas_volume}
\QQRow{酸素 $0.25\,{\rm mol}$}{$5.6\,{\rm L}$}{窒素 $1.5\,{\rm mol}$}{$33.6\,{\rm L}$}
\QQRow{二酸化炭素 $0.40\,{\rm mol}$}{$8.96\,{\rm L}$}{水素 $2.0\,{\rm mol}$}{$44.8\,{\rm L}$}
\QQRow{アンモニア $0.75\,{\rm mol}$}{$16.8\,{\rm L}$}{メタン $1.25\,{\rm mol}$}{$28.0\,{\rm L}$}

\TypeQuestion{標準状態で次の体積を占める気体の物質量を求めなさい。}{gas_volume_to_mol}
\QQRow{水素 $11.2\,{\rm L}$}{$0.5\,{\rm mol}$}{酸素 $2.24\,{\rm L}$}{$0.1\,{\rm mol}$}
\QQRow{窒素 $6.72\,{\rm L}$}{$0.3\,{\rm mol}$}{二酸化炭素 $56.0\,{\rm L}$}{$2.5\,{\rm mol}$}
\QQRow{二酸化炭素 $4.48\,{\rm L}$}{$0.2\,{\rm mol}$}{アンモニア $16.8\,{\rm L}$}{$0.75\,{\rm mol}$}

\AnswerColumnBreak
\LessonQuestionBreak

\TypeQuestion{次の物質量に含まれる粒子の個数を求めなさい。}{mol_to_particles}
\QQRow{酸素分子 $0.25\,{\rm mol}$}{$1.5\times10^{23}$個}{二酸化炭素分子 $1.2\,{\rm mol}$}{$7.2\times10^{23}$個}
\QQRow{ヘリウム原子 $0.050\,{\rm mol}$}{$3.0\times10^{22}$個}{銅原子 $2.5\,{\rm mol}$}{$1.5\times10^{24}$個}
\QQRow{窒素分子 $0.80\,{\rm mol}$}{$4.8\times10^{23}$個}{水素分子 $1.5\,{\rm mol}$}{$9.0\times10^{23}$個}

\TypeQuestion{次の個数の粒子の物質量を求めなさい。}{particles_to_mol}
\QQRow{水素分子 $3.0\times10^{23}$個}{$0.50\,{\rm mol}$}{窒素分子 $1.2\times10^{24}$個}{$2.0\,{\rm mol}$}
\QQRow{ヘリウム原子 $9.0\times10^{22}$個}{$0.15\,{\rm mol}$}{銅原子 $4.8\times10^{23}$個}{$0.80\,{\rm mol}$}
\QQRow{酸素分子 $6.0\times10^{22}$個}{$0.10\,{\rm mol}$}{二酸化炭素分子 $1.5\times10^{23}$個}{$0.25\,{\rm mol}$}

\TypeQuestion{次の質量の気体が標準状態で占める体積を求めなさい。}{mass_to_gas_volume}
\QQRow{水素 $4.0\,{\rm g}$}{$44.8\,{\rm L}$}{酸素 $16\,{\rm g}$}{$11.2\,{\rm L}$}
\QQRow{窒素 $21\,{\rm g}$}{$16.8\,{\rm L}$}{二酸化炭素 $8.8\,{\rm g}$}{$4.48\,{\rm L}$}
\QQRow{アンモニア $17\,{\rm g}$}{$22.4\,{\rm L}$}{メタン $20\,{\rm g}$}{$28\,{\rm L}$}

\TypeQuestion{標準状態で次の体積を占める気体の質量を求めなさい。}{gas_volume_to_mass}
\QQRow{水素 $33.6\,{\rm L}$}{$3\,{\rm g}$}{酸素 $5.60\,{\rm L}$}{$8\,{\rm g}$}
\QQRow{窒素 $44.8\,{\rm L}$}{$56\,{\rm g}$}{二酸化炭素 $2.24\,{\rm L}$}{$4.4\,{\rm g}$}
\QQRow{アンモニア $33.6\,{\rm L}$}{$25.5\,{\rm g}$}{メタン $16.8\,{\rm L}$}{$12\,{\rm g}$}
